\section*{Capítulo \ref{ch:g6:water-underground} --- A água subterrânea}

\begin{solution}{exo:g6:water-underground:1}
\omterm{def:g6:water-underground:permeability}{Permeáveis}: \omterm{def:g3:stones-pebbles-sand:sand}{areia}, \omterm{def:g3:stones-pebbles-sand:pebble}{cascalho}, giz. \omterm{def:g6:water-underground:permeability}{Impermeáveis}: argila, granito sem
rachaduras, folhelho.
\end{solution}

\begin{solution}{exo:g6:water-underground:2}
(a) \omterm{def:g6:water-underground:water-table}{lençol freático}; (b) \omterm{def:g6:water-underground:aquifer}{aquífero}; (c) \omterm{def:g6:water-underground:spring}{fonte}; (d) \omterm{def:g6:water-underground:spring}{poço}.
\end{solution}

\begin{solution}{exo:g6:water-underground:3}
(a) Ela evapora de volta para o ar; ela escorre pela superfície até os
riachos; ela se infiltra no chão (\omterm{def:g6:water-underground:infiltration}{infiltração}). (b) No estacionamento: a
sua superfície dura quase não deixa a água se infiltrar, então a água
escorre. No prado, o solo e a grama deixam boa parte dela se infiltrar.
\end{solution}

\begin{solution}{exo:g6:water-underground:4}
(a) Falso: a argila é muito \omterm{def:g6:water-underground:porosity}{porosa}, mas \omterm{def:g6:water-underground:permeability}{impermeável}, porque os seus \omterm{def:g6:water-underground:porosity}{poros}
são minúsculos. (b) Verdadeiro. (c) Falso: ele segue a forma do terreno, de
modo mais suave, e sobe e desce com as estações. (d) Verdadeiro: cerca de
30~\% de toda a \omterm{def:g2:water-on-earth:salt-fresh}{água doce} é \omterm{def:g6:water-underground:groundwater}{água subterrânea}, contra cerca de um quarto de
um por cento nos \omterm{def:g1:landscapes-around-us:water}{lagos}.
\end{solution}

\begin{solution}{exo:g6:water-underground:5}
$10 \times 25 \div 100 = 2.5$: você consegue despejar 2.5~litros antes que
ele transborde. $10 \times 22 \div 100 = 2.2$: 2.2~litros sairiam pelo furo.
\end{solution}

\begin{solution}{exo:g6:water-underground:6}
(a) $2 \times 1 \times 1 = 2$~metros cúbicos, ou seja, 2000~litros.
(b) $2000 \times 20 \div 100 = 400$~litros. (c) $2000 \times 18 \div 100 =
360$~litros escoam; $400 - 360 = 40$~litros ficam grudados.
\end{solution}

\begin{solution}{exo:g6:water-underground:7}
(a) A \omterm{def:g6:water-underground:spring}{fonte} fica a cerca de 203~m, e o \omterm{def:g1:landscapes-around-us:water}{rio} a cerca de 132~m. (b) A \omterm{def:g6:water-underground:spring}{fonte}
fica cerca de 70~m acima do \omterm{def:g1:landscapes-around-us:water}{rio}: a sua água desce a encosta até o \omterm{def:g1:landscapes-around-us:water}{rio}.
(c) Dentro da encosta de giz, o \omterm{def:g6:water-underground:water-table}{lençol freático} fica abaixo do chão; ele só
encontra a superfície onde a base do giz, apoiada no folhelho \omterm{def:g6:water-underground:permeability}{impermeável},
aparece na encosta.
\end{solution}

\begin{solution}{exo:g6:water-underground:8}
(a) Depois da estação chuvosa: $401 - 336 = 65$~m; depois da estação seca:
$401 - 314 = 87$~m. (b) $336 - 314 = 22$~m.
\end{solution}

\begin{solution}{exo:g6:water-underground:9}
Mais fundo. O \omterm{def:g6:water-underground:water-table}{lençol freático} segue a forma do terreno, mas de modo mais
suave: no alto da \omterm{def:g1:landscapes-around-us:landform}{colina}, ele fica mais alto que embaixo do \omterm{def:g6:water-underground:spring}{poço}~A, mas o
chão sobe muito mais, então a água fica mais longe da superfície (no corte,
a mais de 110~m de profundidade).
\end{solution}

\begin{solution}{exo:g6:water-underground:10}
Os \omterm{def:g6:water-underground:porosity}{poros} da argila são minúsculos: a água fica grudada nas lâminas e não
consegue atravessar a argila até o \omterm{def:g6:water-underground:spring}{poço}. A argila é \omterm{def:g6:water-underground:porosity}{porosa}, mas
\omterm{def:g6:water-underground:permeability}{impermeável}, então só um pouco de água (20~litros de 500) escoa.
\end{solution}

\begin{solution}{exo:g6:water-underground:11}
(a) $30.1 \div 0.26 \approx 116$: mais de cem vezes mais.
(b) $100 - 69 = 31$: sobra cerca de 31~\%, e só a \omterm{def:g6:water-underground:groundwater}{água subterrânea} é cerca de 30 deles.
(c) Ela é de longe a maior reserva de \omterm{def:g2:water-on-earth:salt-fresh}{água doce} líquida; o gelo está
congelado e fica sobretudo longe, nas regiões polares e nas \omterm{def:g1:landscapes-around-us:landform}{montanhas}. E a
\omterm{def:g6:water-underground:groundwater}{água subterrânea} está em quase toda parte embaixo dos nossos pés.
\end{solution}

\begin{solution}{exo:g6:water-underground:12}
Os dois \omterm{def:g6:water-underground:spring}{poços} tiram água do mesmo \omterm{def:g6:water-underground:aquifer}{aquífero} de giz. A bomba potente tirou
água mais depressa do que a \omterm{def:g4:clouds-rain-snow:precipitation}{chuva} reabastece o giz, então o \omterm{def:g6:water-underground:water-table}{lençol freático}
desceu em toda a volta, abaixo do fundo do segundo \omterm{def:g6:water-underground:spring}{poço}, que secou. Os
agricultores poderiam combinar de bombear menos, sobretudo no verão,
dividir um só \omterm{def:g6:water-underground:spring}{poço} ou aprofundar o segundo \omterm{def:g6:water-underground:spring}{poço}; acima de tudo, eles não
deveriam tirar mais água do que a \omterm{def:g4:clouds-rain-snow:precipitation}{chuva} devolve.
\end{solution}

\begin{solution}{pb:g6:water-underground:1}
\textbf{1.} O sítio fica sobre o giz; embaixo dele fica o folhelho. O giz
é \omterm{def:g6:water-underground:permeability}{permeável}, o folhelho não.

\textbf{2.} A \omterm{def:g4:clouds-rain-snow:precipitation}{chuva} se infiltra no giz e para no folhelho. Onde a base do
giz aparece na encosta, o \omterm{def:g6:water-underground:water-table}{lençol freático} encontra o chão e a água sai: é
a \omterm{def:g6:water-underground:spring}{fonte}.

\textbf{3.} $401 - 336 = 65$~m abaixo do chão.

\textbf{4.} O fundo do \omterm{def:g6:water-underground:spring}{poço} fica a $401 - 100 = 301$~m, abaixo do \omterm{def:g6:water-underground:water-table}{lençol
freático}, a 314~m: o \omterm{def:g6:water-underground:spring}{poço} ainda alcança a água, por $314 - 301 =
13$~m.

\textbf{5.} $100 \times 50 \times 2 = 10\,000$~metros cúbicos.

\textbf{6.} $10\,000 \times 20 \div 100 = 2000$~metros cúbicos.

\textbf{7.} $10\,000 \times 18 \div 100 = 1800$~metros cúbicos escoam;
$2000 - 1800 = 200$~metros cúbicos ficam presos.

\textbf{8.} $1800 \times 1000 = 1\,800\,000$~litros.

\textbf{9.} $18 \times 1000 = 18\,000$~litros por dia.

\textbf{10.} $18 \times 7 = 126$~metros cúbicos por semana.

\textbf{11.} A \omterm{def:g4:clouds-rain-snow:precipitation}{chuva} da estação chuvosa se infiltra no giz e o reabastece
todo ano (desde que os sítios não bombeiem mais do que a \omterm{def:g4:clouds-rain-snow:precipitation}{chuva} devolve).

\textbf{12.} $1800 \div 18 = $~\textbf{100 dias}: um pouco mais de três
meses sem nenhuma \omterm{def:g4:clouds-rain-snow:precipitation}{chuva}.
\end{solution}
